% Part: sets-functions-relations
% Chapter: arithmetization
% Section: reflections

\documentclass[../../../include/open-logic-section]{subfiles}

\begin{document}

\olfileid[fa-IR]{sfr}{arith}{ref}
\olsection{چند تأمل فلسفی}

از جزئیات فنی همین بس. اما این جزئیات چه دستاوردی داشتند؟

خب، تقریباً بی‌چون‌وچرا، مقداری ریاضیات محضِ {دل‌انگیز} به ما دادند.
افزون بر این، دستاوردهای مفهومی ژرفی نیز در کار بود. این بینشی ژرف بود که
خاصیت تمامیت تفاوت اساسی میان اعداد حقیقی و گویا را بیان می‌کند. همچنین،
ساخت صریح اعداد حقیقی به‌صورت برش‌های ددکیند، موضوع تحلیل را بر پایه‌ای
استوار قرار می‌دهد. می‌دانیم مفهوم یک \emph{میدان مرتبِ کامل} سازگار است،
زیرا برش‌ها دقیقاً چنین میدانی را تشکیل می‌دهند.

با همهٔ این‌ها، باید چند ملاحظهٔ انتقادی دربارهٔ این دستاوردها مطرح کنیم.

نخست، روشن نیست که اندیشیدن به اعداد حقیقی برحسب برش‌ها از اندیشیدن به
آن‌ها برحسب بسط‌های دهدهیِ آشنا (و شاید نامتناهی) \emph{دقیق‌تر} باشد.
این «ساخت» اخیر اعداد حقیقی تا اندازه‌ای به ساخت اعداد حقیقی به‌کمک دنبالهٔ
کوشی شباهت دارد؛ اما در واقع، ریاضی‌دانان از اوایل سدهٔ هفدهم به بعد اساساً
آن را می‌شناختند (بنگرید به \olref[cauchy]{sec}). افزایش واقعی دقت از این
دریافت آمد که اعداد حقیقی خاصیت تمامیت دارند؛ تواناییِ ساختن اعداد حقیقی
به‌صورت مجموعه‌هایی خاص شاید به‌خودی‌خود چندان جالب نباشد.

حتی کمتر روشن است که حسابی‌سازیِ (بسیار آسان‌ترِ) اعداد صحیح یا گویا، دقت
را در آن حوزه‌ها افزایش دهد. در اینجا شایسته است نکته‌ای ساده را یادآور
شویم. پس از آنکه اعداد صحیح را به‌صورت رده‌های هم‌ارزیِ زوج‌های مرتبِ
اعداد طبیعی \emph{ساختیم}، سپس اعداد گویا را به‌صورت رده‌های هم‌ارزیِ
زوج‌های مرتبِ اعداد صحیح ساختیم، و آنگاه اعداد حقیقی را به‌صورت مجموعه‌هایی
از اعداد گویا ساختیم، بلافاصله این ساخت‌ها را \emph{فراموش می‌کنیم}. به‌ویژه،
هیچ‌کس هرگز نمی‌خواهد در جریان یک برهان ریاضی به این ساخت‌ها \emph{متوسل
شود} (البته به‌استثنای برهانی که نشان دهد ساخت‌ها چنان رفتار می‌کنند که باید).
سخن‌گفتن مستقیم از یک عدد حقیقی بسیار آسان‌تر است تا سخن‌گفتن از مجموعه‌ای
از مجموعه‌هایی از مجموعه‌هایی از مجموعه‌هایی از مجموعه‌هایی از مجموعه‌هایی
از مجموعه‌های اعداد طبیعی.
%(For reals were constructed as sets of rationals; and these were constructed as equivalence classes of ordered pairs of integers, i.e., sets of sets of sets of integers; etc.).
%2/3 = [2,3]
%	i.e. {<2,3>, ... }
%	{ {{2},{2,3}}, ... }
%
%reals are 
%	sets of rationals 
%
%rationals are 
%	sets of sets of sets of integers
%
%integers are
%	sets of sets of sets of naturals
	
از همه تردیدبرانگیزتر این است که این تعریف‌ها به ما بگویند اعداد صحیح،
گویا یا حقیقی \emph{چه هستند}، \emph{به‌لحاظ مابعدالطبیعی}. یعنی محل تردید است که
اعداد حقیقی (مثلاً) \emph{همان} مجموعه‌هایی خاص (از مجموعه‌هایی از
مجموعه‌ها\ldots) باشند. مانع اصلیِ چنین دیدگاهی آن است که ساخت را
می‌شد به شیوه‌های بسیار متفاوتی انجام داد. در مورد اعداد حقیقی، چند ساخت
واقعاً و به‌شکلی جالب متفاوت وجود دارد (بنگرید به \olref[cauchy]{sec}).
اما راهی واقعاً پیش‌پاافتاده برای به‌دست‌آوردن ساخت‌هایی متفاوت این است:
همان‌گونه که در \olref[sfr][rel][ref]{sec} دیدیم، می‌توانستیم زوج‌های مرتب
را اندکی متفاوت تعریف کنیم؛ اگر این مفهوم جایگزین زوج مرتب را به کار
می‌بردیم، ساخت‌های ما دقیقاً به همان خوبی عمل می‌کردند، اما در نهایت به
اشیایی متفاوت می‌رسیدیم. بنابراین ساخت‌های مجموعه‌شناختیِ رقیب بسیاری برای
اعداد صحیح، گویا و حقیقی وجود دارد. اکنون این ادعا که اعداد صحیح (مثلاً)
\emph{این} مجموعه‌ها هستند و نه \emph{آن} مجموعه‌ها، صرفاً دلبخواهانه (و
مایهٔ شرمندگی) بود. (چنان‌که در \olref[sfr][rel][ref]{sec} آمد، این نمونه‌ای
از استدلالی است که \citealt{Benacerraf1965} آن را مشهور کرد.)

نکتهٔ دیگری نیز شایستهٔ طرح است: چیزی بس \emph{عجیب} در ساخت‌های ما وجود
دارد. کار را با اعداد طبیعی آغاز کردیم. سپس اعداد صحیح و «$0$ اعداد صحیح»،
یعنی $ \equivrep{0,0}{\Intequiv}$، را می‌سازیم. اما $0 \neq
\equivrep{0,0}{\Intequiv}$. در واقع، با توجه به ساخت‌های ما، \emph{هیچ}
عدد طبیعی‌ای عدد صحیح نیست. ولی این نتیجه به‌شدت خلاف شهود می‌نماید. افزون
بر این، در \olref[sfr][set][imp]{sec} بی‌آنکه استدلال چندانی عرضه کنیم
ادعا کردیم $\Nat \subseteq \Rat$. اگر این ساخت‌ها دقیقاً به ما بگویند
اعداد \emph{چه هستند}، این ادعا به‌وضوح کاذب بود.

پس اگر از فاصله‌ای دورتر نگاه کنیم، به کجا رسیده‌ایم؟ با کارکردن در یک
نظریهٔ طبیعی مجموعه‌ها و مفروض‌گرفتن اعداد طبیعی، می‌توانیم اعداد صحیح،
گویا و حقیقی را همچون مجموعه‌هایی خاص \emph{در نظر بگیریم}. به این معنا،
می‌توانیم نظریه‌های این موجودیت‌ها را درون یک نظریهٔ مجموعه‌ها \emph{بنشانیم}.
اما اهمیت فلسفی این نشاندن به‌هیچ‌وجه سرراست نیست.

البته هیچ‌یک از این مطالب سخن آخر نیست!{} نکته فقط این است. نشان‌دادن اینکه
حسابی‌سازی اعداد حقیقی \emph{واقعاً} اهمیت فلسفی ژرفی دارد، به استدلال
\emph{فلسفی} بیشتری نیاز خواهد داشت.

%Additionally: for the entire duration of this chapter, we have treated
%the natural numbers as simply \emph{given} to us. But what can we do
%with them? That is the subject matter for the next chapter.

\end{document}
